Would a simpler scorer that cannot be influenced by order do as well? We score each candidate by CLIP prompt--image similarity \citep{radford2021clip,hessel2021clipscore} and return the highest-scoring image. This CLIP baseline has regret 0.130 (95\% CI [0.109, 0.153]) and gain over random 0.069 [0.047, 0.090], compared with Qwen's 0.160 and 0.039; the paired difference in gain favors CLIP by 0.030 [0.005, 0.055]. CLIP also selects a below-mean image less often (28.0\% versus 36.3\%). Yet it raises the stereotype rating relative to random by 0.026 [0.006, 0.046], as much as Qwen does, while sharply reducing missing explicit expectations ($-$0.113). An alignment-oriented scorer therefore does not avoid stereotyped images; this concern is not specific to VLM judges.

Qwen agrees with CLIP on 49.0\% of prompts (chance 27.8\%). Where they agree, Qwen's gain is 0.099 [0.067, 0.130]; where they disagree it is $-$0.018 [$-$0.047, 0.011]. Agreement with an independent, position-invariant scorer thus behaves like order unanimity rather than like agreement with the weaker VLM judge. The CLIP comparison is exploratory (Appendix Table~\ref{tab:baselines}).

Results do not hinge on noisy annotations. Splitting prompts at the median annotator disagreement, Qwen's gain is 0.047 [0.014, 0.079] on high-consensus and 0.033 [0.002, 0.062] on low-consensus prompts (Appendix Table~\ref{tab:robustness}).
